\documentclass[../../main.tex]{subfiles}

\begin{document}

\subsection*{\texttt{fake-reco}}
  Il processo \texttt{fake-reco} attualmente non accetta alcun parametro di
  configurazione (quando sarà aggiunta la possibilità di applicare uno
  \textit{smearing} sarà introdotto un parametro contestuale), non richiede
  alcun dato marcato come Instanced e ritorna un oggetto di tipo
  \texttt{sand::caf::caf\_wrapper} che è Instanced
  (\texttt{sand::caf::caf\_wrapper}, come indicato dal nome, è un
  \textit{wrapper} attorno alla classe \texttt{StandardRecord} per poter
  ereditare da \texttt{ufw::data::base} ed essere usato nel ufw).

  Come già detto il suo scopo attualmente è scrivere la verità Monte Carlo nei
  CAF, sia nei campi a lei destinati, sia in quelli che dovrebbero riportare
  i risultati della ricostruzione.

  Per la lettura dei dati in input \texttt{fake-reco} utilizza sia
  \texttt{genie\_reader} che \texttt{edep\_reader} dato che entrambi hanno
  informazioni esclusive l'uno rispetto all'altro; per esempio la maggior
  parte delle informazioni sul neutrino primario si trovano unicamente su
  genie mentre tutte le informazioni sul tempo o sulle particelle che non
  sono figlie dirette del neutrino sono solo nei dati di EDep-Sim.

\end{document}